\documentclass[12pt,a4paper]{article}

% -------------------------------------------------
% Pacotes
% -------------------------------------------------
\usepackage[utf8]{inputenc}
\usepackage[brazil]{babel}
\usepackage{geometry}
\usepackage{graphicx}
\usepackage{hyperref}
\usepackage{float}
\usepackage{setspace}
\usepackage{xcolor}

\geometry{top=2.5cm,bottom=2.5cm,left=2.5cm,right=2.5cm}
\onehalfspacing

% -------------------------------------------------
% Título
% -------------------------------------------------
\title{
Apontamento de Literatura – Tema 2\\
\large Introdução à Tomada de Decisão Baseada na Evidência e às Revisões Sistemáticas
}
\author{
Disciplina: Revisão Sistemática e Meta-análises\\
Pós-graduação (Mestrado e Doutoramento)
}
\date{}

\begin{document}

\maketitle

% =================================================
\section{Evidência: conceito fundamental}
% =================================================
O termo \textbf{evidência} refere-se a qualquer informação empírica ou teórica utilizada para sustentar uma afirmação, inferência ou decisão. No contexto científico, evidência não é sinónimo de verdade absoluta, mas sim de \textbf{conhecimento provisório}, sujeito a revisão, crítica e falseamento.

\textbf{Características da evidência científica:}
\begin{itemize}
    \item Obtida por métodos explícitos e sistemáticos;
    \item Reprodutível ou verificável;
    \item Sujeita à revisão por pares;
    \item Possui incerteza quantificável.
\end{itemize}

\begin{figure}[H]
    \centering
    \includegraphics[width=0.6\textwidth]{fig1_evidencia.png}
    \caption{Diagrama conceptual da evidência científica}
    \label{fig:evidencia}
\end{figure}

% =================================================
\section{Evidência científica}
% =================================================
A \textbf{evidência científica} resulta da aplicação do método científico a fenómenos naturais ou sociais. Distingue-se de:
\begin{itemize}
    \item Evidência anedótica;
    \item Opiniões de especialistas não sistematizadas;
    \item Experiência pessoal não controlada.
\end{itemize}

Em termos epistemológicos, a evidência científica é:
\begin{itemize}
    \item \textbf{Empírica} (baseada em dados observáveis);
    \item \textbf{Teórica} (interpretada à luz de modelos e hipóteses);
    \item \textbf{Probabilística} (nunca absolutamente certa).
\end{itemize}

% =================================================
\section{Práticas Baseadas na Evidência (EBP)}
% =================================================
As práticas baseadas na evidência consistem na integração sistemática de:
\begin{enumerate}
    \item Melhor evidência científica disponível;
    \item Experiência profissional especializada;
    \item Contexto, valores e necessidades dos utilizadores finais.
\end{enumerate}

\begin{figure}[H]
    \centering
    \includegraphics[width=0.5\textwidth]{fig2_ebp.png}
    \caption{Triângulo da prática baseada na evidência (EBP)}
    \label{fig:ebp}
\end{figure}

% =================================================
\section{Necessidade de falseabilidade}
% =================================================
Segundo Karl Popper, uma afirmação científica deve ser \textbf{falseável}, isto é, deve ser possível conceber uma observação ou experimento que a contradiga.

\begin{figure}[H]
    \centering
    \includegraphics[width=0.7\textwidth]{fig3_falseabilidade.png}
    \caption{Fluxo de hipótese, teste e falseamento}
    \label{fig:falseabilidade}
\end{figure}

% =================================================
\section{Fluxo da evidência: da geração à implementação}
% =================================================
O ciclo da evidência segue as etapas: geração, síntese, transferência, implementação e avaliação.

\begin{figure}[H]
    \centering
    \includegraphics[width=0.8\textwidth]{fig4_fluxo_evidencia.png}
    \caption{Fluxo da evidência desde a geração até à implementação}
    \label{fig:fluxo_evidencia}
\end{figure}

% =================================================
\section{O ecossistema da evidência}
% =================================================
Envolve múltiplos atores: investigadores, revisores, instituições, sociedade civil, decisores.

\begin{figure}[H]
    \centering
    \includegraphics[width=0.6\textwidth]{fig5_ecossistema.png}
    \caption{Rede do ecossistema da evidência}
    \label{fig:ecossistema}
\end{figure}

% =================================================
\section{Síntese da evidência}
% =================================================
A síntese da evidência é necessária para reduzir incertezas e evitar decisões baseadas em estudos isolados.

\begin{figure}[H]
    \centering
    \includegraphics[width=0.5\textwidth]{fig6_sintese.png}
    \caption{Pirâmide da síntese da evidência}
    \label{fig:sintese}
\end{figure}

% =================================================
\section{Tipos de revisão}
% =================================================
\begin{itemize}
    \item \textbf{Revisão narrativa}: não sistemática, subjetiva, contextualização teórica.
    \item \textbf{Revisão de escopo}: mapeamento de área e lacunas.
    \item \textbf{Revisão rápida}: metodologia simplificada, decisões urgentes.
    \item \textbf{Revisão sistemática}: metodologia explícita, avaliação crítica, síntese qualitativa e quantitativa.
\end{itemize}

\begin{figure}[H]
    \centering
    \includegraphics[width=0.6\textwidth]{fig7_tipos_revisao.png}
    \caption{Comparação entre tipos de revisão}
    \label{fig:tipos_revisao}
\end{figure}

% =================================================
\section{Revisão sistemática}
% =================================================
Principais elementos:
\begin{itemize}
    \item Questão estruturada (PICO);
    \item Protocolo pré-definido;
    \item Critérios de inclusão/exclusão;
    \item Busca abrangente;
    \item Avaliação da qualidade;
    \item Síntese transparente.
\end{itemize}

\begin{figure}[H]
    \centering
    \includegraphics[width=0.8\textwidth]{fig8_fluxo_revisao.png}
    \caption{Fluxograma do processo de revisão sistemática}
    \label{fig:fluxo_revisao}
\end{figure}

% =================================================
\section{Preparação para uma revisão sistemática}
% =================================================
Checklist:
\begin{itemize}
    \item Verificar revisões existentes;
    \item Avaliar viabilidade;
    \item Definir objetivos e escopo;
    \item Escolher diretrizes metodológicas;
    \item Registo no PROSPERO.
\end{itemize}

\begin{figure}[H]
    \centering
    \includegraphics[width=0.6\textwidth]{fig9_checklist.png}
    \caption{Checklist de preparação para revisão sistemática}
    \label{fig:checklist}
\end{figure}

% =================================================
\section{Processo completo de revisão sistemática}
% =================================================
O processo completo envolve: questão → protocolo → registo → busca → seleção → extração → avaliação → síntese → reporte.

\begin{figure}[H]
    \centering
    \includegraphics[width=0.9\textwidth]{fig10_processo_revisao.png}
    \caption{Fluxo completo da revisão sistemática}
    \label{fig:processo_revisao}
\end{figure}

% =================================================
\section{Considerações finais}
% =================================================
Para mestrados e doutoramentos, a revisão sistemática é:
\begin{itemize}
    \item Uma postura epistemológica;
    \item Um instrumento de decisão científica;
    \item Um pilar da ciência moderna baseada na evidência.
\end{itemize}

\end{document}
